Our \robot\ can jump to the edge or to the centre of the frame (which delimits
the drawing area on the computer screen), without changing
its orientation and without drawing anything. It is exactly as if we could
 carry it
by hand and put it down at the edge or at the centre of the frame, as one can do with a
remote-controlled car. The computer does it for us thanks to the commands
given in the table below and illustrated in figure
\ref{fig:ecrangdhbc} (page \pageref{fig:ecrangdhbc}) (the underlined letters give the abbreviation of the names of the commands):
\begin{table}[htbp]
\begin{center}
\begin{tabular}{@{}|l|p{11cm}|@{}} \hline
{\tt gl()} & (\underline{g}o to the \underline{l}eft) places the \robot\ on the left of the screen, without changing its
                  vertical position or its orientation, and without drawing anything.\\ \hline
{\tt gr()} & (\underline{g}o to the \underline{r}ight) places the \robot\ on the right of the screen, without changing its
                  vertical position or its orientation, and without drawing anything.\\ \hline
{\tt gt()} & (\underline{g}o to the \underline{t}op) places the \robot\ at the top of the screen, without changing its
                  horizontal position or its orientation, and without drawing anything.\\ \hline
{\tt gb()} & (\underline{g}o to the \underline{b}ottom) places the \robot\ at the bottom of the screen, without changing its
                  horizontal position or its orientation, and without drawing anything.\\ \hline
{\tt gc()} & (\underline{g}o to the \underline{c}entre) places the \robot\ at the centre of the screen, without changing its orientation, and
                  without drawing anything.\\ \hline
\end{tabular}
\end{center}
%\label{tab:ag-ad-eh-eb-ce}
\caption{Jumps of the \protect\robot\ to the edge and to the centre of the frame}
\end{table}

{\samepage %1
\begin{figure}[ht]
\includevspacePostScript[7.2cm]{ecran-ag-ad-eh-eb-ce}{0.95}{0.95}
\caption{Jumps of the \protect\robot\ to the edge and to the centre of the frame}
\label{fig:ecrangdhbc}
\end{figure}
\needspace{6\baselineskip}\noindent \rule{137mm}{0.1mm}\par\nopagebreak %---------------------------------------------
\begin{exercice}{
Write a \python\ program that draws crosses at the edge and at the centre of the screen
as shown in the figure below.
\noindent $\Box$}\end{exercice}
\vspace*{15mm}
\noindent\begin{minipage}{\textwidth}\LEcaptype{figure}
\includevspacePostScript[6.5cm]{CroixDeCoins}{0.95}{0.95}
 \caption{Crosses at the edge and at the centre of the frame}
 %\label{fig:CroixDeCoins}
\end{minipage}
} %1


\newpage
\begin{corrige}{\ \\

\vspace{-\baselineskip}
\input python/CornerCrosses.tex
\noindent $\Box$}\end{corrige}

%%%%%%%%%%%%%%%%%%%%%%%%%%%%%%%%%%%%%%%%%%%%%%%%%%%%%%%%%%%%%%%%%%%%%%%%%%%
\noindent \framebox[\textwidth]{
\begin{minipage}{\myparboxwidth}
\setlength{\parindent}{2em}
\small \sf
\vspace{\abovedisplayskip}
\section*{Addenda}

The commands that move the \robot\ to the edge of the frame will be defined in \python\ later, after
we have explained the notions of test and iteration. We introduce them right now so that we can easily draw
friezes.  In the next chapter these friezes with repetitive patterns will allow us to introduce simply
the notion of bounded iteration.

The command {\idt gc()}, which brings the \robot\ back to the centre, will be essential as soon as we use a Cartesian coordinate system to designate the points
of the grid.

\vspace{\abovedisplayskip}
\end{minipage}
}
%%%%%%%%%%%%%%%%%%%%%%%%%%%%%%%%%%%%%%%%%%%%%%%%%%%%%%%%%%%%%%%%%%%%%%%%%%%%
\newpage
\thispagestyle{empty}